%ExpressionSemantics.tex


\ksec{Expressions Semantics (KerML 8.4.4.9)}\label{sec:expressions}

%\comment{Unfortunately, an  \texttt{Expression} when declared with \texttt{expr} is a \texttt{Step}--not calculation of a value.
%Other parts of Expression define what we normally think of as expression.
%Other than \texttt{Expression} being a step (feature) there seems to be no difference of \texttt{expr} from a \texttt{function} other than general confusion because an expression is something used in assignment (and other things) to calculate a value.
%}
%
%No timing semantics.

\pn An \bemph{expression} returns a \bemph{value} (\ref{sec:value})--not performs a \texttt{Step}.  Therefore, \texttt{Expression} evaluation takes no time, and thus is suitable for binding, initialization, defaults, and constraints.   If modeling calculation of a value during operation is desired, use  \texttt{Performances::Computation}.  

\pn This requires changing the specialization of metaclass \texttt{Expression} from \texttt{Step} to \texttt{Feature}.
\ecaption{metaclass Expression}
\begin{lstlisting}
  metaclass Expression specializes Feature {
  	derived var feature parameter : Feature[0..*] subsets input;
  	derived var feature 'function' : Function[0..1];
  	derived var feature result : Feature[1..1] subsets output, parameter;
  }
\end{lstlisting}

 \pn Expressions (and Functions and Predicates) are \emph{pure}.
 Expressions have no side effects.\footnote{to have formal semantics}

\pn Expressions do not reference \texttt{self}.

\pn Expressions may have 0 or more inputs.

\pn Expressions produce a single result, whose value is uniquely determined by the values of their inputs.\footnote{to have formal semantics}

\pn Specializations of Expressions may not change parameter identifiers, nor their order,
and may add parameters only at the end (thus preserving parameter order of their generalizations).

\pn Type inferencing of expressions are contravariant (see \ref{sec:contra})

\begin{bnf}
Expression =
  FeaturePrefix
  'expr' FeatureDeclaration ValuePart?
  FunctionBody 
\end{bnf}

\subsection{Extending Expression for SFS}\label{sec:extendingexpression}

\pn Currently, SFS are added to KerML and SysML models with ``Domain'' language annotations.
Hopefully, these annotations will use the same expression as KerML and SysML.
By adding non-executable expressions, invariant clauses (\texttt{inv}) could express semantics directly.

\pn Add non-executable expressions to \texttt{BaseExpression}.
A \texttt{NonExecutableExpression} cannot be used for initialization, assignment, or transition conditions.
Validations after parsing could easily detect misuse of non-executable expressions.

\begin{bnf}
BaseExpression ::=
  NullExpression
  | LiteralExpression
  | FeatureReferenceExpression
  | MetadataAccessExpression
  | InvocationExpression
  | ConstructorExpression
  | BodyExpression
  | {\color{orange}NonExecutableExpression}
\end{bnf}

\begin{bnf}
Param ::= Identifier ( ',' Identifier )* '~' Type
\end{bnf}

\begin{bnf}
NonExecutableExpression ::=
  'timed' PrimaryExpression 'at' PrimaryExpression
| 'shifted' PrimaryExpression 'by' PrimaryExpression
| 'forall' Param ( ',' Param )* ( 'in' RangeOrBoolean )?
    'are' BooleanExpression
| 'forall' Param ( ',' Param )* ( 'in' RangeOrBoolean )?
    'forall' Param ( ',' Param )* ( 'in' RangeOrBoolean )?
    'are' BooleanExpression
| 'exists' Param ( ',' Param )* ( 'in' RangeOrBoolean )?
    'that' BooleanExpression
| 'numberof'
    Identifier ( ',' Identifier )* '~' Type ( 'in' RangeOrBoolean )?
    'that' BooleanExpression
| ( 'productof' | 'sumof' )
    Identifier ( ',' Identifier )* '~' Type ( 'in' RangeOrBoolean )?
    'that' PrimaryExpression
\end{bnf}
\begin{description}
\item[\texttt{Param}:] a single \texttt{forall}/\texttt{exists} may bind several
  differently-typed \texttt{Param} groups at once (e.g.\
  \texttt{forall x\textasciitilde{}Occurrence, y,z\textasciitilde{}Region are \ldots}).  The second
  \texttt{forall \ldots\ forall \ldots\ are} form (two \texttt{forall} keywords, not comma-joined)
  is needed for formulas like \texttt{LFU}/\texttt{LIN}/\texttt{EXPNS}/\texttt{APAR}.
\item[\texttt{numberof}/\texttt{productof}/\texttt{sumof}:] unlike \texttt{forall}/\texttt{exists},
  these three bind only a single shared type across their identifier list -- no \texttt{Param}
  grouping.
\item[\texttt{PrimaryExpression}:] \texttt{timed}/\texttt{shifted}'s operands are not separately
  typed as time-valued or integer-valued at the grammar level; that distinction, if enforced at
  all, happens elsewhere.  \texttt{shifted \ldots\ by \ldots} parses but is not yet semantically
  implemented.
\item[\texttt{BooleanExpression}:] expression resulting in \emph{true} or  \emph{false}.
\end{description}

\begin{bnf}
RangeOrBoolean ::=
  PrimaryExpression RangeSymbol PrimaryExpression
  | BooleanExpression
\end{bnf}

\begin{bnf}
RangeSymbol ::= '..' | '.,' | ',.' | ',,'
\end{bnf}



